%%%%%%%%%%%%%%%%%%%%%%%%%%%%%%%%%%%%%%%%%%%%%%%%%%%%%%%%%%%%%%%%%%%%%%%%%%%%%%%%
%%%%%%%%%%%%%%%%%%%%%%%%%%%%%%%%%%%%%%%%%%%%%%%%%%%%%%%%%%%%%%%%%%%%%%%%%%%%%%%%
%%%%%%%%%%%%%%%%%%% BASE DOCUMENT SETTINGS
%%%%%%%%%%%%%%%%%%%%%%%%%%%%%%%%%%%%%%%%%%%%%%%%%%%%%%%%%%%%%%%%%%%%%%%%%%%%%%%%
%%%%%%%%%%%%%%%%%%%%%%%%%%%%%%%%%%%%%%%%%%%%%%%%%%%%%%%%%%%%%%%%%%%%%%%%%%%%%%%%
$if(pdf-x)$
%% PDF/X metadata, which pdfx reads from \jobname.xmpdata
\begin{filecontents}[force]{\jobname.xmpdata}
$if(title-meta)$\Title{$title-meta$}$endif$
$if(author-meta)$\Author{$author-meta$}$endif$
$if(lang)$\Language{$lang$}$endif$
$if(publisher)$\Publisher{$for(publisher)$$publisher$$endfor$}$endif$
$if(rights)$\Copyright{$rights$}$endif$
\end{filecontents}
$endif$

%% Set document class
\documentclass%
[11pt]
{book} % Set document Class

%% Set page size and margins
\usepackage{geometry}
  \geometry{
  paperheight=$if(paper-height)$$paper-height$$else$228.6mm$endif$,  %actual height measurement of the page (default is 6x9 inches)
  paperwidth=$if(paper-width)$$paper-width$$else$152.4mm$endif$,    %actual width measurement of the page
  $if(copyedit)$
    includeall, %include header, footer, and margin notes inside of text area calculation (in copyedit mode, show margin notes off to the side)
  $else$
    includehead, %include header inside of text area calculation (footer and margin notes aren't used in this template)
  $endif$
  top=$if(top-margin)$$top-margin$$else$30.0mm$endif$, % top margin
  bottom=$if(bottom-margin)$$bottom-margin$$else$27.0mm$endif$, %bottom margin
  inner=$if(inner-margin)$$inner-margin$$else$22.0mm$endif$, %inner/gutter margin
  outer=$if(outer-margin)$$outer-margin$$else$19.0mm$endif$, %outer margin
  $if(copyedit)$
    showframe %show lines drawn around text, footer, header, and margin note areas
  $endif$
  }

%% Set line spacing
% 1.0 = single space
% 1.3 = 1.5 space
% 1.6 = double space
\linespread{1.3}

% Useful for switching line spacing for small sections in sub-documents
\usepackage{setspace}
  
%% Set no additional space between paragraphs
\setlength{\parskip}{0pt}

%% Use this package to enable tables  
\usepackage{longtable,booktabs}

%% Hide the toprule
\renewcommand{\toprule}{}

%% Hide the bottomrule
\renewcommand{\bottomrule}{}

\setlength{\LTpost}{0pt}

%% Set language and hyphenation rules from lang (e.g., fr-CA is canadien)
\usepackage[$if(babel-lang)$$babel-lang$$else$english$endif$]{babel}

%% The word used for "Chapter" in chapter headings and running headers
$if(chapter-name)$
\setlocalecaption{\languagename}{chapter}{$chapter-name$}
$endif$

%% The word used for "Part" on part pages
$if(part-name)$
\setlocalecaption{\languagename}{part}{$part-name$}
$endif$

%% Hyphenation exceptions: words that should not be split across lines, or with explicit break points
%% (no-hyphenate-list is kept for backward compatibility)
\hyphenation{$no-hyphenate-list$ $hyphenation-exceptions$}

%% Set encoding
%\usepackage[utf8]{inputenc}

%% Enable widow/orphan control
\usepackage[all,defaultlines=3]{nowidow}

%% Prevent lines from spilling over
\setlength{\emergencystretch}{3em}
\providecommand{\tightlist}{\setlength{\itemsep}{0pt}\setlength{\parskip}{0pt}}

%%%%%%%%% Print boxes, IF USING xelatex (PDF/X itself is set by pdf-x below) %%%%%%%%%
% Thanks to LSinev for most of this;
% https://gist.github.com/LSinev/74ea05423c9208d6677d

% TrimBox is the page size in bp, calculated from paper-width and paper-height
% (rounded to avoid exceeding the page size)
\newcommand{\trimboxwidth}{\fpeval{round(\paperwidth*72/72.27, 2)}}
\newcommand{\trimboxheight}{\fpeval{round(\paperheight*72/72.27, 2)}}
\usepackage{atbegshi}
\AtBeginShipout{% %A hook that is executed for every page (after first one)
    \special{pdf: put @thispage
      <<
        /TrimBox [0 0 \trimboxwidth\space\trimboxheight]
      >>
    }
}
\special{pdf: put @thispage
  <<
    /TrimBox [0 0 \trimboxwidth\space\trimboxheight]
  >>
}
\special{pdf:put @catalog
<<
  /PageMode /UseNone
>>
}%

\usepackage{datetime} % for \pdfdate command

$if(pdf-x)$
%% Export PDF as PDF/X-1a:2001
%% pdfx loads hyperref and loads xcolor with CMYK colors, so their options are passed here
\PassOptionsToPackage{bookmarks=false}{hyperref}
\PassOptionsToPackage{dvipsnames,svgnames}{xcolor}
\usepackage[x-1a1]{pdfx}
%% pdfx insets the TrimBox from the page edges, but the page is the trim size
\makeatletter
\gdef\pdfx@everypage@xetex{/TrimBox [0 0 \trimboxwidth\space\trimboxheight]}
\makeatother
\special{pdf:put @pages << /TrimBox [0 0 \trimboxwidth\space\trimboxheight] >>}
$endif$

%% URL / Hyperlink style
\usepackage[bookmarks=false]{hyperref}
\NoHyper


\hypersetup{pdfstartpage={}}

\hypersetup{unicode=true,
$if(title-meta)$
            pdftitle={$title-meta$}, %set pdftitle metadata from Yaml front matter
$endif$
            pdfinfo={% Setting some more PDF/X stuff for xelatex
                ModDate={D:\pdfdate},% PDF/X document should have a modification date
                Trapped={False},% PDF/X document should have Trapped tag set
            },
$if(author-meta)$
            pdfauthor={$author-meta$}, %set pdfauthor metadata from Yaml front matter
$endif$
$if(keywords)$
            pdfkeywords={$for(keywords)$$keywords$$sep$; $endfor$}, %set pdfkeywords metadata from Yaml front matter
$endif$
$if(colorlinks)$
            colorlinks=true,
            linkcolor=$if(linkcolor)$$linkcolor$$else$Maroon$endif$,
            citecolor=$if(citecolor)$$citecolor$$else$Blue$endif$,
            urlcolor=$if(urlcolor)$$urlcolor$$else$Blue$endif$,
$else$
            pdfborder={0 0 0}, %turn off the url border that normally gets added to urls in pdfs
$endif$
            breaklinks=true}
            
\urlstyle{same}  % don't use monospace font for urls

$for(header-includes)$
$header-includes$
$endfor$

%%%%%%%%%%%%%%%%%%%%%%%%%%%%%%%%%%%%%%%%%%%%%%%%%%%%%%%%%%%%%%%%%%%%%%%%%%%%%%%%
%%%%%%%%%%%%%%%%%%%%%%%%%%%%%%%%%%%%%%%%%%%%%%%%%%%%%%%%%%%%%%%%%%%%%%%%%%%%%%%%
%%%%%%%%%%%%%%%%%%% HEADERS AND FOOTERS
%%%%%%%%%%%%%%%%%%%%%%%%%%%%%%%%%%%%%%%%%%%%%%%%%%%%%%%%%%%%%%%%%%%%%%%%%%%%%%%%
%%%%%%%%%%%%%%%%%%%%%%%%%%%%%%%%%%%%%%%%%%%%%%%%%%%%%%%%%%%%%%%%%%%%%%%%%%%%%%%%
%% Needed for headers and footers
\usepackage{fancyhdr}
%% Useful for stacking text in a title
\usepackage{stackengine}
\def\stackalignment{l}

%% Prevent headers from appearing on empty pages
\usepackage{emptypage}

%% Page number placement (header, bottom-center, or bottom-outer)
\def\pagenumberposition{$if(page-number-position)$$page-number-position$$else$header$endif$}
\def\bottomcenter{bottom-center}
\def\bottomouter{bottom-outer}

%% Remove the header from the first page of every chapter (bottom page numbers are kept)
\fancypagestyle{plain}{
\fancyhf{} %remove/clear header/footer content
\renewcommand{\headrulewidth}{0.0pt} %remove black line/"rule" from header
\ifx\pagenumberposition\bottomcenter \fancyfoot[C]{\primaryfont \thepage}\fi
\ifx\pagenumberposition\bottomouter \fancyfoot[RO,LE]{\primaryfont \thepage}\fi
}

%% Foreword and preface pages: roman page numbers, with "Foreword" or "Preface" in the header
\makeatletter
\fancypagestyle{frontmatter}{
\fancyhf{}
\renewcommand{\headrulewidth}{0.0pt}
\ifx\pagenumberposition\bottomcenter
  \fancyfoot[C]{\primaryfont \thepage}
\else\ifx\pagenumberposition\bottomouter
  \fancyfoot[RO,LE]{\primaryfont \thepage}
\else
  \fancyhead[RO,LE]{\primaryfont\bfseries \thepage}
\fi\fi
\fancyhead[LO]{\primaryfont\bfseries \textsc \leftmark}
\fancyhead[RE]{\primaryfont\bfseries \textsc \@title}
}
\makeatother

%% Set the header height in: in, mm, cm, pt
\setlength{\headheight}{0.10in}

%% Set the space between header and text
\setlength{\headsep}{.25in}

%%%%%%%%%%%%%%%%%%%%%%%%%%%%%%%%%%%%%%%%%%%%%%%%%%%%%%%%%%%%%%%%%%%%%%%%%%%%%%%%
%%%%%%%%%%%%%%%%%%%%%%%%%%%%%%%%%%%%%%%%%%%%%%%%%%%%%%%%%%%%%%%%%%%%%%%%%%%%%%%%
%%%%%%%%%%%%%%%%%%% FONTS
%%%%%%%%%%%%%%%%%%%%%%%%%%%%%%%%%%%%%%%%%%%%%%%%%%%%%%%%%%%%%%%%%%%%%%%%%%%%%%%%
%%%%%%%%%%%%%%%%%%%%%%%%%%%%%%%%%%%%%%%%%%%%%%%%%%%%%%%%%%%%%%%%%%%%%%%%%%%%%%%%
%% Use TrueType system fonts
\usepackage{fontspec}

%% Curly quotes and dashes (must come before fonts are defined)
\defaultfontfeatures{Ligatures=TeX}

%% For editing font names
\usepackage{xstring}

%% Used to set some things to all caps
\usepackage{textcase}

%% Provides symbols for flourishes (scene dividers, chapter title decorations)
\usepackage{adforn}

%% Fontawesome icons for various flourishes
\usepackage{fontawesome5}

%% Provides Creative Commons Icons
\usepackage{ccicons}

%% Color macros for fonts
\usepackage[dvipsnames,svgnames]{xcolor}

%% TikZ (for hand-drawn marker highlights and ornaments)
\usepackage{tikz}
\usetikzlibrary{backgrounds,calc,decorations.pathmorphing}

%% Vector ornaments for the backmatter pages
\usepackage{pgfornament}

%% Add support for strikethrough (\st{} command)
\usepackage{soul}

%% Provides support for drop caps
\usepackage{lettrine}

%% Support for custom font sizes
\usepackage{anyfontsize}

\def\romance{romance}
\def\scifi{scifi}
\def\oldwest{oldwest}
\def\blockletter{blockletter}

%% Drop cap color (black unless drop-cap-color is specified)
\def\dropcapcolor{$if(drop-cap-color)$$drop-cap-color$$else$black$endif$}
\renewcommand{\LettrineFontHook}{\color{\dropcapcolor}}

$if(drop-cap-style)$

  %% If you want the first letter to be a calligraphy font
  \ifx\$drop-cap-style$\romance
    \newfontfamily\dropcapfont{texgyrechorus-mediumitalic.otf}
    \renewcommand{\LettrineFontHook}{\dropcapfont\color{\dropcapcolor}}
    \renewcommand*{\DefaultLhang}{.25}
    \renewcommand*{\DefaultLoversize}{.7}
  \fi

  %% A science-fiction look
  \ifx\$drop-cap-style$\scifi
    \newfontfamily\dropcapfont{Orbitron-Bold.ttf}
    \renewcommand{\LettrineFontHook}{\dropcapfont\color{\dropcapcolor}}
    \renewcommand*{\DefaultLoversize}{.5}
  \fi

  %% Old Western look
  \ifx\$drop-cap-style$\oldwest
    \newfontfamily\dropcapfont{ChunkFive-Regular.ttf}
    \renewcommand{\LettrineFontHook}{\dropcapfont\color{\dropcapcolor}}
    \renewcommand*{\DefaultLoversize}{.5}
  \fi

  %% Or a block letter
  \ifx\$drop-cap-style$\blockletter
    \usepackage{Romantik}
    \renewcommand{\LettrineFontHook}{\usefont{U}{Romantik}{xl}{n}\color{\dropcapcolor}}
    \renewcommand*{\DefaultLoversize}{.5}
    \setlength{\DefaultFindent}{0.2em}
    \setlength{\DefaultNindent}{0.2em}
  \fi

$endif$

%% Set Main Fonts (Cochineal, unless main-font is specified)
$if(main-font)$
\def\primaryfontfile{$main-font$}
\def\primaryfontfeatures{$if(main-font-bold)$BoldFont = $main-font-bold$,$endif$$if(main-font-italic)$ItalicFont = $main-font-italic$,$endif$$if(main-font-bold-italic)$BoldItalicFont = $main-font-bold-italic$$endif$}
\StrDel{$main-font$}{.ttf}[\primaryfontname]
\StrDel{\primaryfontname}{.otf}[\primaryfontname]
$else$
% Cochineal is an extended version of Crimson with true small caps
\def\primaryfontfile{Cochineal-Roman.otf}
\def\primaryfontfeatures{BoldFont = Cochineal-Bold.otf, ItalicFont = Cochineal-Italic.otf, BoldItalicFont = Cochineal-BoldItalic.otf}
\def\primaryfontname{Cochineal}
$endif$

% Runs a fontspec command (e.g., \setmainfont) with the primary font and its bold/italic variants
\newcommand\withprimaryfont[1]{%
  \edef\primaryfontcommand{\unexpanded{#1}{\primaryfontfile}[\primaryfontfeatures]}%
  \primaryfontcommand}

\withprimaryfont{\setmainfont}
\withprimaryfont{\newfontfamily\sansfont}
\withprimaryfont{\newfontfamily\primaryfont}
\withprimaryfont{\newfontfamily\chapterheadingfont}

$if(author-font)$
  \newfontfamily\authorfont{$author-font$}
$endif$

$if(title-font)$
  \newfontfamily\titlefont{$title-font$}
$else$
  \withprimaryfont{\newfontfamily\titlefont}
$endif$

%%%%%%%%%%%%%%%%%%%%%%%%%%%%%%%%%%%%%%%%%%%%%%%%%%%%%%%%%%%%%%%%%%%%%%%%%%%%%%%%
%%%%%%%%%%%%%%%%%%%%%%%%%%%%%%%%%%%%%%%%%%%%%%%%%%%%%%%%%%%%%%%%%%%%%%%%%%%%%%%%
%%%%%%%%%%%%%%%%%%% METADATA and MACROS
%%%%%%%%%%%%%%%%%%%%%%%%%%%%%%%%%%%%%%%%%%%%%%%%%%%%%%%%%%%%%%%%%%%%%%%%%%%%%%%%
%%%%%%%%%%%%%%%%%%%%%%%%%%%%%%%%%%%%%%%%%%%%%%%%%%%%%%%%%%%%%%%%%%%%%%%%%%%%%%%%

%% Set Title from Yaml Metadata
\title{$title$}

%% Set subtitle if it exists in Yaml Metadata
$if(subtitle)$
\newcommand{\subtitle}{$subtitle$}
$endif$

%% Set author from Yaml Metadata
\author{$for(author)$$author$$endfor$}

%% Set editor from Yaml Metadata
\def\editor{$for(contributors)$$contributors.editor$$endfor$}

%% Set cover artist/illustrator from Yaml Metadata
\def\coverartist{$for(contributors)$$contributors.cover-artist$$endfor$}
\def\illustrator{$for(contributors)$$contributors.illustrator$$endfor$}

%% Set designer from Yaml Metadata
\def\designer{$for(contributors)$$contributors.designer$$endfor$}

%% Set publisher from Yaml Metadata
\def\publisher{$for(publisher)$$publisher$$endfor$}

%% Set website from Yaml Metadata
\def\website{$for(website)$$website$$endfor$}

%% Use the year instead of the full date
\date{$year$}

%% Converts number to string
\usepackage{fmtcount}
% load the language's number words before the first chapter heading, since loading them inside a heading breaks its line
\AtBeginDocument{\sbox0{\Numberstringnum{1}}}

%% Chapter numbers in words ("Chapter One") or digits
\newif\ifchapternumberwords
$if(chapter-number-words)$\chapternumberwordstrue$endif$
\newcommand{\chapternumber}{\ifchapternumberwords\Numberstringnum{\thechapter}\else\thechapter\fi}
\newcommand{\CHAPTERNUMBER}{\ifchapternumberwords\NUMBERstringnum{\thechapter}\else\thechapter\fi}
%% Part numbers follow the same choice, but use Roman numerals instead of digits
\newcommand{\partnumber}{\ifchapternumberwords\Numberstringnum{\arabic{part}}\else\Roman{part}\fi}

%%%%%%%%%%%%%%%%%%%%%%%%%%%%%%%%%%%%%%%%%%%%%%%%%%%%%%%%%%%%%%%%%%%%%%%%%%%%%%%%
%%%%%%%%%%%%%%%%%%%%%%%%%%%%%%%%%%%%%%%%%%%%%%%%%%%%%%%%%%%%%%%%%%%%%%%%%%%%%%%%
%%%%%%%%%%%%%%%%%%% TITLE SETTINGS
%%%%%%%%%%%%%%%%%%%%%%%%%%%%%%%%%%%%%%%%%%%%%%%%%%%%%%%%%%%%%%%%%%%%%%%%%%%%%%%%
%%%%%%%%%%%%%%%%%%%%%%%%%%%%%%%%%%%%%%%%%%%%%%%%%%%%%%%%%%%%%%%%%%%%%%%%%%%%%%%%
%% Needed for styling chapter title headings
\usepackage{titlesec}
\usepackage{fp}

%% Hand-drawn "marker" highlight: a single-line highlight with wavy, marker-like edges
\newcommand{\marker}[1]{%
  \if\relax\detokenize{#1}\relax\else
  \tikz[baseline=(X.base)]{%
    \node[inner sep=2pt, outer sep=0pt] (X) {#1};%
    \begin{scope}[on background layer]%
      \fill[yellow!80, decorate, decoration={random steps, segment length=4pt, amplitude=1.2pt}]%
        ([xshift=-2pt,yshift=-2pt]X.south west) rectangle ([xshift=2pt]$$(X.north east)!0.40!(X.south east)$$);%
    \end{scope}}%
  \fi}

% justification
\usepackage{ragged2e}
\usepackage{quoting}

% https://tex.stackexchange.com/questions/27977/centered-but-left-aligned-paragraph
\newenvironment{raggedrighttitle}[1][3cm]
  {\begin{quoting}[leftmargin=#1,rightmargin=#1]\RaggedRight}
  {\end{quoting}}

% Turn off subsection numbering for titles
\renewcommand{\thesection}{}

%% Style the chapter title that appears on the first page of each chapter
\titlespacing*{\chapter}
{0pt}%left margin
{150pt}%top margin
{25pt}%bottom margin

% This tells it to style the title format for chapters (or use a style from fncychap, if requested)
$if(fancy-chapter)$
  \usepackage[$fancy-chapter$]{fncychap}
$else$
  \input{chapter-heading.tex}
$endif$

% Unnumbered chapters (e.g., a prologue or epilogue) don't increment the running header,
% so set it to the chapter's title (rather than leaving the previous chapter's title there)
\makeatletter
\let\fp@chapter\chapter
\renewcommand{\chapter}{\fp@chapterimage\@ifstar{\fp@starchapter}{\fp@chapter}}
\newcommand{\fp@starchapter}[1]{\fp@chapter*{#1}\markboth{#1}{}}

%% Artwork above each chapter heading in the main matter.
%% It is drawn in the space above the heading, so it works with every chapter heading theme (but not fncychap).
$if(chapter-image)$
$if(fancy-chapter)$
\newcommand{\fp@chapterimage}{}
$else$
\usepackage{eso-pic}
\newcommand{\fp@chapterimage}{%
  \if@mainmatter
    % start the chapter's page first, so that the image is added to it
    \if@openright\cleardoublepage\else\clearpage\fi
    \AddToShipoutPictureFG*{\AtTextUpperLeft{\put(\LenToUnit{0.5\textwidth},\LenToUnit{-75pt}){%
      \makebox(0,0){\includegraphics[width=$if(chapter-image-width)$$chapter-image-width$$else$1.5in$endif$,height=120pt,keepaspectratio]{$chapter-image$}}}}}%
  \fi}
$endif$
$else$
\newcommand{\fp@chapterimage}{}
$endif$

%% Part pages for every chapter heading theme: "Part One" in spaced small caps, a short rule, and the part's title
\newcommand{\fp@parttitle}[1]{\if\relax\detokenize{#1}\relax\else\rule{2cm}{0.4pt}\endgraf\vspace{6mm}\Huge #1\fi}
\titleformat{\part}[display]
  {\filcenter\chapterheadingfont}%
  {\Large\scshape\addfontfeature{LetterSpace=8}\partname\ \partnumber}%
  {6mm}%
  {\fp@parttitle}
\assignpagestyle{\part}{empty}
\makeatother

% Set spacing between chapters in TOC to be closer
\usepackage[titles]{tocloft}
\setlength{\cftbeforechapskip}{2pt}
\renewcommand{\cftpartfont}{\normalfont\scshape}
\renewcommand{\cftpartpagefont}{\normalfont}
\setlength{\cftbeforepartskip}{12pt}

% Add spacing for the chapter title that appears on the first page of each chapter
%%%%%%%%%%%%%%%%%%%

%%%%%%%%%%%%%%%%%%%%%%%%%%%%%%%%%%%%%%%%%%%%%%%%%%%%%%%%%%%%%%%%%%%%%%%%%%%%%%%%
%%%%%%%%%%%%%%%%%%%%%%%%%%%%%%%%%%%%%%%%%%%%%%%%%%%%%%%%%%%%%%%%%%%%%%%%%%%%%%%%
%%%%%%%%%%%%%%%%%%% START THE DOCUMENT
%%%%%%%%%%%%%%%%%%%%%%%%%%%%%%%%%%%%%%%%%%%%%%%%%%%%%%%%%%%%%%%%%%%%%%%%%%%%%%%%
%%%%%%%%%%%%%%%%%%%%%%%%%%%%%%%%%%%%%%%%%%%%%%%%%%%%%%%%%%%%%%%%%%%%%%%%%%%%%%%%
\begin{document}

%% Start with the frontmatter
% These pages don't count for document page numbering
\frontmatter

%%%%%%%%%%%%%%%%%%%%%%%%%%%
%%%%%%%%%%%%%%%%%%%%%%%%%%%
%% Create a praise page
%%%%%%%%%%%%%%%%%%%%%%%%%%%
%%%%%%%%%%%%%%%%%%%%%%%%%%%
$if(praise)$
  \input{praise.tex}
  \clearpage{\thispagestyle{empty}\cleardoublepage}
$endif$

%%%%%%%%%%%%%%%%%%%%%%%%%%%
%%%%%%%%%%%%%%%%%%%%%%%%%%%
%% Create a half titlepage
%%%%%%%%%%%%%%%%%%%%%%%%%%%
%%%%%%%%%%%%%%%%%%%%%%%%%%%

\input{half-titlepage.tex}

% End the page
\clearpage

%%%%%%%%%%%%%%%%%%%%%%%%%%%
%%%%%%%%%%%%%%%%%%%%%%%%%%%
%% Create the series page
%%%%%%%%%%%%%%%%%%%%%%%%%%%
%%%%%%%%%%%%%%%%%%%%%%%%%%%

$if(series-book1)$
  \input{seriespage.tex}
$endif$

\clearpage{\thispagestyle{empty}\cleardoublepage}

%%%%%%%%%%%%%%%%%%%%%%%%%%%
%%%%%%%%%%%%%%%%%%%%%%%%%%%
%% Create a titlepage
%%%%%%%%%%%%%%%%%%%%%%%%%%%
%%%%%%%%%%%%%%%%%%%%%%%%%%%
% No headers/footers
\pagestyle{empty}
% Center and single space everything
  \begin{spacing}{.75}
  \begin{center}

% Add a vertical space
  \vspace*{$if(title-page-top-margin)$$title-page-top-margin$$else$4cm$endif$}

% Write the Book Title
\def\titlefontsize{$if(title-font-pointsize)$$title-font-pointsize$$else$50$endif$}
\FPeval\titlefontbaseline{\titlefontsize*1.2}
$if(title-latex)$
  \makeatletter {\titlefont\fontsize{\titlefontsize}{\titlefontbaseline}\selectfont $title-latex$} \par \makeatother
$else$
  \makeatletter {\titlefont\fontsize{\titlefontsize}{\titlefontbaseline}\selectfont \textbf{\textsc{\@title}}} \par \makeatother
$endif$

% Add a vertical space
  \vspace*{.5cm}
    
% Write the subtitle
$if(subtitle-latex)$
  \Large \primaryfont $subtitle-latex$ \par
$elseif(subtitle)$
  \Large \primaryfont $subtitle$ \par
$endif$
    
% Add a vertical space and switch to regular title page font
  \vspace*{4cm}
  \normalfont \sansfont \normalsize

% Write the author name
$if(author-by-header)$
  $author-by-header$ \par
$endif$

$if(author-font)$
  {\Large \authorfont \makeatletter \MakeUppercase\@author} \makeatother \par
$else$
  {\Large \makeatletter \MakeUppercase\@author} \makeatother \par
$endif$

% Illustrator
$for(contributors)$$if(contributors.illustrator)$ \vspace*{4cm} \normalsize Illustrated by \par \Large \illustrator \par$endif$$endfor$

% Stop centering and single spacing everything
  \end{center}
  \end{spacing}

% End the page
\clearpage

%%%%%%%%%%%%%%%%%%%%%%%%%%%
%%%%%%%%%%%%%%%%%%%%%%%%%%%
%% Create a copyright page
%%%%%%%%%%%%%%%%%%%%%%%%%%%
%%%%%%%%%%%%%%%%%%%%%%%%%%%
% No headers/footers

\input{copyright.tex}

\clearpage

%%%%%%%%%%%%%%%%%%%%%%%%%%%
%%%%%%%%%%%%%%%%%%%%%%%%%%%
%% Create a dedication page
%%%%%%%%%%%%%%%%%%%%%%%%%%%
%%%%%%%%%%%%%%%%%%%%%%%%%%%
$if(dedication)$
    % No headers/footers
    \pagestyle{empty}

    % Center everything
      \begin{center}

    % Vertically center
    \topskip0pt
    \vspace*{\fill}
  
    % Write the Dedication
      \normalfont\normalsize $dedication$ \par

    % Stop centering
    \vspace*{\fill}
        \end{center}
    
    % End the page
    \clearpage
$endif$

%%%%%%%%%%%%%%%%%%%%%%%%%%%
%%%%%%%%%%%%%%%%%%%%%%%%%%%
%% Create an epigraph page
%%%%%%%%%%%%%%%%%%%%%%%%%%%
%%%%%%%%%%%%%%%%%%%%%%%%%%%
$if(epigraph)$
    \pagestyle{empty}
    \vspace*{\stretch{1}}
    \begin{center}
    \begin{minipage}{0.7\textwidth}
      \setlength{\parindent}{0pt}
      \setlength{\parskip}{6pt}
      % a large opening quote hangs in the margin beside the first line, its top level with the capitals
      \def\epigraphquote{$epigraph-quote$}\def\epigraphnoquote{false}%
      \ifx\epigraphquote\epigraphnoquote\else
      \setbox0=\hbox{{\normalfont\fontsize{48pt}{48pt}\selectfont\color{black!60}“}}%
      \makebox[0pt][r]{\smash{\raisebox{\dimexpr\fontcharht\font`T-\ht0\relax}{\copy0}}\hspace{3pt}}%
      \fi
      {\itshape $epigraph$\par}
    $if(epigraph-source)$
      \vspace{8pt}
      {\raggedleft\textemdash\,\textsc{$epigraph-source$}\par}
    $endif$
    \end{minipage}
    \end{center}
    \vspace*{\stretch{2}}
    \clearpage
$endif$

%%%%%%%%%%%%%%%%%%%%%%%%%%%
%%%%%%%%%%%%%%%%%%%%%%%%%%%
%% Create a table of contents
%%%%%%%%%%%%%%%%%%%%%%%%%%%
%%%%%%%%%%%%%%%%%%%%%%%%%%%
$if(toc)$
    $if(toc-title)$
      \renewcommand*\contentsname{$toc-title$}
    $endif$
    $if(beamer)$
      \begin{frame}[allowframebreaks]
      $if(toc-title)$
        \frametitle{$toc-title$}
      $endif$
      \tableofcontents[hideallsubsections]
      \end{frame}
    $else$
      {
      $if(colorlinks)$
        \hypersetup{linkcolor=$if(toccolor)$$toccolor$$else$$endif$}
      $endif$
      \setcounter{tocdepth}{$toc-depth$}
      \tableofcontents
      }
    $endif$
$endif$

%%%%%%%%%%%%%%%%%%%%%%%%%%%
%%%%%%%%%%%%%%%%%%%%%%%%%%%
%% Create a foreword
%%%%%%%%%%%%%%%%%%%%%%%%%%%
%%%%%%%%%%%%%%%%%%%%%%%%%%%
$if(foreword)$
\input{foreword.tex}

\clearpage{\thispagestyle{empty}\cleardoublepage}
$endif$

%%%%%%%%%%%%%%%%%%%%%%%%%%%
%%%%%%%%%%%%%%%%%%%%%%%%%%%
%% Create a preface
%%%%%%%%%%%%%%%%%%%%%%%%%%%
%%%%%%%%%%%%%%%%%%%%%%%%%%%
$if(preface)$
\input{preface.tex}

\clearpage{\thispagestyle{empty}\cleardoublepage}
$endif$

%%%%%%%%%%%%%%%%%%%%%%%%%%%
%%%%%%%%%%%%%%%%%%%%%%%%%%%
%% Create the chapters
%%%%%%%%%%%%%%%%%%%%%%%%%%%
%%%%%%%%%%%%%%%%%%%%%%%%%%%
% This styles the header/footer for normal 'chapter' pages
\makeatletter
\pagestyle{fancy}
%\renewcommand{\chaptermark}[1]{\markboth{#1}{}} % This removes the 'Chapter x' from the header
% Chapter with no title shows only 'Chapter N', without a trailing ' -- '
\renewcommand{\chaptermark}[1]{\markboth{\if\relax\detokenize{#1}\relax\@chapapp\ \chapternumber\else\@chapapp\ \chapternumber\ --\ {#1}\fi}{}}

\fancyhf{} % clear default header/footer content
\ifx\pagenumberposition\bottomcenter
  \fancyfoot[C]{\primaryfont \thepage}
\else\ifx\pagenumberposition\bottomouter
  \fancyfoot[RO,LE]{\primaryfont \thepage}
\else
  \fancyhead[RO,LE]{\primaryfont\bfseries \thepage} % add page numbers to the header
\fi\fi
\fancyhead[LO]{\primaryfont\bfseries \textsc \leftmark} % add subtitle to header
\fancyhead[RE]{\primaryfont\bfseries \textsc \@title}
\renewcommand{\headrulewidth}{0.0pt} 
\makeatother

% This says to start the page numbering
\mainmatter

% Write the body/chapters
\primaryfont \normalsize \flushbottom $body$

%%%%%%%%%%%%%%%%%%%%%%%%%%%
%%%%%%%%%%%%%%%%%%%%%%%%%%%
%% Create the backmatter
%%%%%%%%%%%%%%%%%%%%%%%%%%%
%%%%%%%%%%%%%%%%%%%%%%%%%%%

% add blank even page before backmatter starts
\clearpage{\thispagestyle{empty}\cleardoublepage}

% Label as backmatter
\backmatter

% No headers/footers
\pagestyle{empty}

% Add space before/after the page title
\titlespacing*{\section}
{0pt}%left
{150pt}%top
{25pt}%bottom

% This designs the page title
\titleformat{\section}[display]{\bfseries\chapterheadingfont}{}{0.5mm}{\centering \Large} % Don't say 'chapter X'

% Don't indent new paragraphs
\setlength{\parindent}{0mm}

% Ornament rule under a backmatter page title
\newcommand{\backmatterornament}{%
  {\centering\rule[0.5ex]{2cm}{0.4pt}\hspace{0.8em}{\large\adforn{21}}\hspace{0.8em}\rule[0.5ex]{2cm}{0.4pt}\par}%
  \vspace{20pt}}

% Ornament that closes a backmatter page
\newcommand{\backmatterend}{\par\vspace{18pt}{\centering\large\adforn{36}\par}}

% Backmatter page title: letter-spaced capitals with an ornament rule beneath
\newcommand{\backmattertitle}[1]{%
  \vspace*{120pt}%
  {\centering\chapterheadingfont\bfseries\Large\addfontfeature{LetterSpace=10}\MakeTextUppercase{#1}\par}%
  \vspace{12pt}%
  \backmatterornament}

% Backmatter title set flush left over a heavy short rule, like a letter's heading
\newcommand{\backmatterlefttitle}[1]{%
  \vspace*{100pt}%
  {\raggedright\chapterheadingfont\bfseries\LARGE\addfontfeature{LetterSpace=8}\MakeTextUppercase{#1}\par}%
  \vspace{8pt}%
  {\raggedright\rule{4em}{3pt}\par}%
  \vspace{36pt}}

% Flourish that closes the acknowledgments and excerpt (a pgfornament number, or none)
\def\backmatterflourishnumber{$if(backmatter-flourish)$$backmatter-flourish$$else$88$endif$}
\def\flourishnone{none}
\newcommand{\backmatterflourish}{\ifx\backmatterflourishnumber\flourishnone\else\pgfornament[width=4cm]{\backmatterflourishnumber}\fi}

% Flourish above the afterword's sign-off (a pgfornament number, or none)
\def\afterwordflourishnumber{$if(afterword-flourish)$$afterword-flourish$$else$none$endif$}
\newcommand{\afterwordflourish}{\ifx\afterwordflourishnumber\flourishnone\else\pgfornament[width=3cm]{\afterwordflourishnumber}\par\vspace{6pt}\fi}

% Teaser page for a book preview: small-caps kicker, the book's title, and its author, centered on the page
\newcommand{\backmatterteaser}[3]{%
  \vspace*{\stretch{1}}%
  {\centering{\primaryfont\scshape\addfontfeature{LetterSpace=8}\MakeTextLowercase{#1}}\par
   \vspace{16pt}%
   {\titlefont\Huge #2}\par
   \vspace{20pt}%
   {\scshape\addfontfeature{LetterSpace=5}\MakeTextLowercase{#3}}\par}}

%%%%%%%%%%%%%%%%%%%%%%%%%%%
%%%%%%%%%%%%%%%%%%%%%%%%%%%
%% Create the afterword, acknowledgments, and reading group guide
%%%%%%%%%%%%%%%%%%%%%%%%%%%
%%%%%%%%%%%%%%%%%%%%%%%%%%%

$if(afterword)$
\input{afterword.tex}

\clearpage{\thispagestyle{empty}\cleardoublepage}
$endif$

$if(acknowledgments)$
\input{acknowledgments.tex}

\clearpage{\thispagestyle{empty}\cleardoublepage}
$endif$

$if(reading-group-guide)$
\input{reading-group-guide.tex}

\clearpage{\thispagestyle{empty}\cleardoublepage}
$endif$

%%%%%%%%%%%%%%%%%%%%%%%%%%%
%%%%%%%%%%%%%%%%%%%%%%%%%%%
%% Create the bio page
%%%%%%%%%%%%%%%%%%%%%%%%%%%
%%%%%%%%%%%%%%%%%%%%%%%%%%%

$if(author-bio)$
\input{bio.tex}

\clearpage{\thispagestyle{empty}\cleardoublepage}
$endif$

%%%%%%%%%%%%%%%%%%%%%%%%%%%
%%%%%%%%%%%%%%%%%%%%%%%%%%%
%% Create the excerpt of the next book
%%%%%%%%%%%%%%%%%%%%%%%%%%%
%%%%%%%%%%%%%%%%%%%%%%%%%%%

$if(excerpt)$
\input{excerpt.tex}

\clearpage{\thispagestyle{empty}\cleardoublepage}
$endif$

%%%%%%%%%%%%%%%%%%%%%%%%%%%
%%%%%%%%%%%%%%%%%%%%%%%%%%%
%% Create the colophon (last page of the book)
%%%%%%%%%%%%%%%%%%%%%%%%%%%
%%%%%%%%%%%%%%%%%%%%%%%%%%%

$if(colophon)$
\input{colophon.tex}

\clearpage
$endif$

% Include any additional content from the metadata
$include-after$

% End the document
\end{document} 
 